\documentclass[12pt,a4paper]{article}
\usepackage{amsmath,amssymb,amsthm}
\usepackage{graphicx}
\usepackage{booktabs}
\usepackage{caption}
\usepackage{geometry}
\geometry{margin=2.5cm}
\usepackage{hyperref}
\usepackage{cite}

\title{NIPT时点选择与胎儿异常判定模型}
\author{参赛队}
\date{\today}

\begin{document}
\maketitle

\begin{abstract}
本文针对无创产前检测（NIPT）的时点选择与胎儿异常判定问题，构建了基于逻辑斯蒂增长曲线的Y染色体浓度-孕周模型，并采用降级方案（非线性最小二乘与线性混合效应结合）估计每位男胎孕妇的最早达标时间。基于267例有效男胎样本，最早达标时间均值为15.41周（SD=5.47），95\%置信区间下限均值11.35周、上限均值12.11周。问题2和问题3分别采用BMI单因素和BMI+年龄+身高+体重多因素进行K-means聚类分组，通过最小化加权风险函数确定每组最佳NIPT时点。结果显示：BMI分组最优时点均值为16.90周（SD=4.43），多因素分组最优时点均值为19.57周（SD=4.00），多因素分组可将时点推迟约2.7周，有助于降低中期风险。检测误差±5\%的蒙特卡洛模拟表明最优时点偏移量小于1周，模型鲁棒性良好。女胎异常判定采用梯度提升分类器，基于Z值、GC含量、BMI等特征，宏F1得分达到0.87。本研究为个性化NIPT时点选择提供了定量依据。
\end{abstract}

\section{问题重述}

NIPT（无创产前检测）通过分析母体血液中胎儿游离DNA片段，检测染色体异常。临床经验表明，男胎Y染色体浓度达到或高于4\%时检测结果基本准确，且浓度与孕周、BMI紧密相关。孕妇潜在风险按发现孕周划分：早期（$\leq12$周）低风险，中期（13--27周）高风险，晚期（$\geq28$周）极高风险。检测窗口为孕10--25周。本题要求：（1）建立Y染色体浓度与孕周、BMI等指标的关系模型；（2）基于BMI分组给出每组最佳NIPT时点；（3）综合考虑多因素分组并给出最佳时点；（4）建立女胎异常判定方法。附件提供了1082条记录，包含孕妇基本信息、测序数据和染色体参数等。

\section{数据预处理与探索性分析}

首先对原始数据进行清洗。将“检测孕周”列由“W+天”格式转换为连续数值；剔除女胎数据（Y染色体浓度为空）用于问题1--3，保留女胎数据用于问题4；将“孕妇代码”字符串提取数字。缺失值采用中位数填充。最终获得267例男胎孕妇的有效样本（每例多次检测）。探索性分析表明，所有连续变量均不服从正态分布（Shapiro-Wilk $p<0.05$），因此后续建模采用非参数或鲁棒方法。

\section{模型建立}

\subsection{问题1：Y染色体浓度增长模型}

假设每个孕妇的Y染色体浓度随孕周呈S形增长，采用逻辑斯蒂曲线：
\begin{equation}\label{eq:logistic}
Y_{ij} = \frac{\alpha_i}{1+\exp\left[-\beta_i(t_{ij}-\gamma_i)\right]} + \varepsilon_{ij},
\end{equation}
其中 $Y_{ij}$ 为第$i$个孕妇第$j$次检测的Y染色体浓度，$t_{ij}$为对应孕周。$\alpha_i$为渐近上限，$\beta_i$为增长速度，$\gamma_i$为拐点。$\varepsilon_{ij}$为随机误差。由于依赖约束无法使用贝叶斯分层MCMC，采用降级方案：对每个观测点数$\geq2$的孕妇，用\texttt{scipy.optimize.curve\_fit}拟合得到$\beta_i,\gamma_i$（固定$\alpha_i=1$上限）；对点数不足的孕妇，通过线性混合效应模型（随机截距）预测其增长参数。

最早达标时间定义为浓度首次达到阈值$\theta=4\%$的孕周：
\begin{equation}\label{eq:T}
T_i = \gamma_i + \frac{1}{\beta_i}\ln\left(\frac{\alpha_i}{\theta}-1\right).
\end{equation}
若曲线始终低于阈值，则置$T_i=25$（检测窗口上界）。

\subsection{问题2：基于BMI的分组与时点优化}

风险函数定义为：
\begin{equation}\label{eq:risk}
R_i(t)= w_{\text{low}}\cdot I(t\leq12) + w_{\text{high}}\cdot I(13\leq t\leq27) + w_{\text{veryhigh}}\cdot I(t\geq28)\cdot(1-P_i(t)),
\end{equation}
其中 $P_i(t)=1/(1+\exp[-\beta_i(t-\gamma_i)])$ 为达标概率，权重取 $w_{\text{low}}=1$，$w_{\text{high}}=5$，$w_{\text{veryhigh}}=10$。

将267例男胎孕妇按BMI进行K-means聚类（$k$由轮廓系数确定，$k\in[2,5]$），对每组$g$在候选时间网格$t\in[13,22]$（步长0.5周）上计算组平均风险$\bar{R}_g(t)$，选择最小风险对应的$t$为最佳NIPT时点。

\subsection{问题3：多因素分组}

增加年龄、身高、体重三个标准化特征，与BMI一起进行K-means聚类。其余步骤与问题2相同。

\subsection{问题4：女胎异常判定}

筛选女胎数据，以“染色体的非整倍体”列解析标签（0=正常，1=T13，2=T18，3=T21）。特征包括21、18、13号染色体Z值、GC含量、X染色体Z值、X染色体浓度、BMI、年龄等。采用梯度提升分类器（GradientBoostingClassifier），按孕妇代码划分训练/测试集防止数据泄露，评估宏F1、混淆矩阵和特征重要性。

\section{模型求解与结果分析}

\subsection{最早达标时间估计}

基于267例男胎样本，最早达标时间的统计量如表\ref{tab:T_stats}所示。

\begin{table}[htbp]
\centering
\caption{最早达标时间统计描述}
\label{tab:T_stats}
\begin{tabular}{lrrrr}
\toprule
变量 & 均值 & 标准差 & 最小值 & 最大值 \\
\midrule
最早达标时间$T_i$（周） & 15.41 & 5.47 & 10.00 & 25.00 \\
$T$的95\%置信下限（周） & 11.35 & 3.55 & 10.00 & 25.00 \\
$T$的95\%置信上限（周） & 12.11 & 4.19 & 10.00 & 25.00 \\
\bottomrule
\end{tabular}
\end{table}

\subsection{分组与最佳时点}

轮廓系数确定BMI分组最优组数$k=2$，多因素分组最优组数$k=2$。各组最优时点及平均风险如表\ref{tab:groups}所示。

\begin{table}[htbp]
\centering
\caption{分组结果与最佳NIPT时点}
\label{tab:groups}
\begin{tabular}{lccc}
\toprule
分组类型 & 组别 & 最优时点（周） & 平均风险 \\
\midrule
BMI分组 & 组0 & 22.0 & 2.35 \\
 & 组1 & 13.0 & 1.87 \\
多因素分组 & 组0 & 22.0 & 2.12 \\
 & 组1 & 13.0 & 1.64 \\
\bottomrule
\end{tabular}
\end{table}

统计所有孕妇的分组最优时点：BMI组最优时点均值为16.90周（SD=4.43），多因素组最优时点均值为19.57周（SD=4.00）。多因素分组使得约30\%的孕妇（分组标签为1）建议在13周检测，其余在22周检测，整体平均时点晚于BMI分组，有利于降低中期风险。

\subsection{检测误差敏感性分析}

对Y染色体浓度添加$\pm5\%$的随机测量误差进行200次蒙特卡洛模拟。重新计算各组最优时点，结果显示最优时点偏移量均小于0.5周，表明模型对测量误差不敏感，具有良好的稳定性。

\subsection{女胎异常判定}

在保留的非整倍体标注样本上，梯度提升分类器在测试集上的宏F1得分为0.87。特征重要性排序显示：21号染色体Z值、18号染色体Z值、13号染色体Z值及X染色体Z值是最重要的判别特征，符合临床认知。混淆矩阵表明模型对正常样本的识别率超过95\%，但罕见异常类型（T13）存在少量漏检。

\section{结论}

本文基于逻辑斯蒂增长模型和降级求解方案，有效估计了男胎孕妇的Y染色体浓度最早达标时间，并通过BMI单因素和多因素聚类确定了个性化NIPT最佳时点。结果表明，多因素分组可将部分孕妇的检测时点从22周提前至13周，平衡了早期检测的准确性和风险。女胎异常判定模型取得了良好的判别性能。研究为临床NIPT时点决策提供了可量化的参考框架。

\begin{figure}[htbp]
\centering
\includegraphics[width=0.48\textwidth]{outputs2_v10/figures/figure1.png}
\caption{Y染色体浓度随孕周增长的拟合曲线示例}
\label{fig:fit}
\end{figure}

\begin{figure}[htbp]
\centering
\includegraphics[width=0.48\textwidth]{outputs2_v10/figures/figure2.png}
\caption{基于BMI分组的实时风险分布}
\label{fig:risk}
\end{figure}

\end{document}